\chapter{Objetivos}

\section{Objetivo Principal}

Ampliar el sistema de gestión de deportistas desarrollado en el proyecto anterior, incorporando algoritmos de ordenamiento, búsqueda y selección basados en la estrategia de \textit{Divide y Vencerás}, con el propósito de analizar su comportamiento teórico y experimental sobre distintos tamaños de entrada y aplicarlos a funcionalidades concretas del sistema, tales como ordenamiento por criterios, búsqueda de registros, generación de rankings, obtención del $k$-ésimo mejor deportista y consulta de deportistas dentro de un rango de puntaje. \\

\section{Objetivos Secundarios}

\begin{itemize}
    \item Implementar algoritmos de ordenamiento basados en \textit{Divide y Vencerás}, específicamente \textit{Merge Sort}, \textit{Merge-Insertion Sort} y \textit{Quick Sort}, permitiendo ordenar deportistas según distintos criterios como ID, puntaje, competencias, nombre y equipo. \\

    \item Implementar distintas variantes de \textit{Quick Sort} utilizando partición de Lomuto, considerando la selección del pivote como último elemento, primer elemento, elemento aleatorio y mediana de tres, con el fin de evaluar cómo esta decisión afecta el rendimiento del algoritmo.  \\

    \item Implementar algoritmos de búsqueda sobre arreglos ordenados, incluyendo búsqueda binaria recursiva, búsqueda binaria por rangos, búsqueda exponencial y búsqueda por interpolación, aplicándolos sobre los datos del sistema de deportistas.  \\

    \item Implementar el algoritmo \textit{Quick Select} para obtener el $k$-ésimo mejor deportista según un criterio determinado, especialmente el puntaje, evitando ordenar completamente el arreglo cuando solo se requiere una posición específica.  \\

    \item Integrar los algoritmos desarrollados al sistema existente, manteniendo la generación, carga y almacenamiento de registros en archivos CSV, así como las opciones de menú y ejecución por línea de comandos.  \\

    \item Diseñar y ejecutar \textit{benchmarks} experimentales que permitan medir los tiempos de ejecución de los algoritmos implementados en distintos tamaños de entrada y en escenarios de mejor caso, caso promedio y  peor caso.  \\

    \item Comparar los resultados experimentales obtenidos con el análisis teórico de complejidad temporal, identificando coincidencias, diferencias prácticas, efectos de las variantes implementadas y posibles limitaciones de la medición.  \\
\end{itemize}